% Part: first-order-logic
% Chapter: tableaux
% Section: identity

\documentclass[../../../include/open-logic-section]{subfiles}

\begin{document}

\olfileid[id]{fol}{tab}{ide}

\olsection{\usetoken{P}{tableau} dengan \usetoken{S}{identity}}

!!^{tableau}s dengan !!{identity} memerlukan aturan inferensi tambahan.
Aturan-aturan untuk $\eq$ adalah sebagai berikut ($t$, $t_1$, dan $t_2$
merupakan suku tertutup):

\begin{defish}
\AxiomC{}
\RightLabel{$\eq$}
\UnaryInfC{\sFmla{\True}{\eq[t][t]}}
\DisplayProof
\hfill
\AxiomC{\sFmla{\True}{\eq[t_1][t_2]}}
\noLine
\UnaryInfC{\sFmla{\True}{!A(t_1)}}
\RightLabel{$\TRule{\True}{\eq}$}
\UnaryInfC{\sFmla{\True}{!A(t_2)}}
\DisplayProof
\hfill
\AxiomC{\sFmla{\True}{\eq[t_1][t_2]}}
\noLine
\UnaryInfC{\sFmla{\False}{!A(t_1)}}
\RightLabel{$\TRule{\False}{\eq}$}
\UnaryInfC{\sFmla{\False}{!A(t_2)}}
\DisplayProof
\end{defish}
Perhatikan bahwa, berbeda dari semua aturan lainnya,
$\TRule{\True}{\eq}$ dan $\TRule{\False}{\eq}$ mengharuskan
\emph{dua} !!{formula} bertanda telah muncul pada cabang tersebut, yakni
$\sFmla{\True}{\eq[t_1][t_2]}$ dan $\sFmla{S}{!A(t_1)}$.

\begin{ex}
Jika $s$ dan $t$ merupakan suku tertutup, maka
$\eq[s][t], !A(s) \Proves !A(t)$:
\begin{oltableau}
  [\sFmla{\False}{\formula{A}(t)}, just = \TAss
    [\sFmla{\True}{\eq[s][t]}, just = \TAss
      [\sFmla{\True}{\formula{A}(s)}, just = \TAss
        [\sFmla{\True}{\formula{A}(t)}, just={\TRule{\True}{\eq}[2, 3]}, close]
      ]
    ]
  ]
\end{oltableau}
Hal ini mungkin telah dikenal sebagai prinsip substitutabilitas hal-hal
yang identik, atau Hukum Leibniz.

!!^{tableau}s juga membuktikan bahwa $\eq$ bersifat simetris, yakni
$\eq[s_1][s_2] \Proves \eq[s_2][s_1]$:
\begin{oltableau}
  [\sFmla{\False}{\eq[s_2][s_1]}, just = \TAss
    [\sFmla{\True}{\eq[s_1][s_2]}, just = \TAss
      [\sFmla{\True}{\eq[s_1][s_1]}, just = {$\eq$}
        [\sFmla{\True}{\eq[s_2][s_1]}, just = {\TRule{\True}{\eq}[2, 3]}, close]
      ]
    ]
  ]
\end{oltableau}
Di sini, baris~$2$ merupakan !!{formula} prasyarat pertama
$\sFmla{\True}{\eq[s_1][s_2]}$ untuk $\TRule{\True}{\eq}$. Baris~$3$
merupakan prasyarat kedua, yang berbentuk
$\sFmla{\True}{!A(s_1)}$---anggap $!A(x)$ sebagai $\eq[x][s_1]$; dengan
demikian, $!A(s_1)$ adalah $\eq[s_1][s_1]$ dan $!A(s_2)$ adalah
$\eq[s_2][s_1]$.

Tableau juga membuktikan bahwa $\eq$ bersifat transitif, yakni
$\eq[s_1][s_2], \eq[s_2][s_3] \Proves \eq[s_1][s_3]$:
\begin{oltableau}
  [\sFmla{\False}{\eq[s_1][s_3]}, just = \TAss
    [\sFmla{\True}{\eq[s_1][s_2]}, just = \TAss
      [\sFmla{\True}{\eq[s_2][s_3]}, just = \TAss
        [\sFmla{\True}{\eq[s_1][s_3]}, just = {\TRule{\True}{\eq}[3, 2]}, close]
      ]
    ]
  ]
\end{oltableau}
Dalam !!{tableau} ini, !!{formula} prasyarat pertama untuk
$\TRule{\True}{\eq}$ adalah baris~$3$,
$\sFmla{\True}{\eq[s_2][s_3]}$ ($s_2$ berperan sebagai~$t_1$ dan $s_3$
sebagai~$t_2$). Prasyarat kedua, yang berbentuk
$\sFmla{\True}{!A(s_2)}$, adalah baris~$2$. Di sini, anggap $!A(x)$
sebagai $\eq[s_1][x]$; dengan demikian, $!A(s_2)$ adalah
$\eq[s_1][s_2]$ (yakni baris~$2$) dan $!A(s_3)$ adalah
!!{formula}~$\eq[s_1][s_3]$ dalam kesimpulan.
\end{ex}

\begin{prob}
Berikan !!{tableau} tertutup untuk hal-hal berikut:
\begin{enumerate}
\item $\sFmla{\False}{\lforall[x][\lforall[y][((x = y \land !A(x))
      \lif !A(y))]]}$
\item $\sFmla{\False}{\lexists[x][(!A(x) \land
    \lforall[y][(!A(y) \lif y = x)])]}$,\\
  $\sFmla{\True}{\lexists[x][!A(x)] \land
    \lforall[y][\lforall[z][((!A(y) \land !A(z)) \lif y = z)]]}$
\end{enumerate}
\end{prob}

\end{document}
